\chapter{Trích Đoạn Mã Nguồn Cốt Lõi}
\label{app:codelistings}

Phụ lục này trích dẫn các đoạn mã nguồn nhúng cốt lõi trong firmware của trạm quan trắc chất lượng không khí, làm minh chứng kỹ thuật cho các phân tích thiết kế ở Chương 4.

\section{Khởi tạo và Tạo các Task FreeRTOS Pinned to Core (main.cpp)}
Đoạn mã cấu hình các tác vụ song song chạy trên lõi kép của ESP32 trong hàm \texttt{setup()}:

\begin{lstlisting}[language=C++]
// Create TaskSensors running on Core 1 (Priority 2)
xTaskCreatePinnedToCore(
    TaskSensorsCode,
    "TaskSensors",
    4096,
    NULL,
    2,                  // Priority 2
    &xTaskSensorsHandle,
    1                   // Pinned to Core 1
);

// Create TaskOLED running on Core 1 (Priority 1)
xTaskCreatePinnedToCore(
    TaskOLEDCode,
    "TaskOLED",
    4096,
    NULL,
    1,                  // Priority 1
    &xTaskOLEDHandle,
    1                   // Pinned to Core 1
);

// Create TaskNetwork running on Core 0 (Priority 1)
xTaskCreatePinnedToCore(
    TaskNetworkCode,
    "TaskNetwork",
    8192,
    NULL,
    1,                  // Priority 1
    &xTaskNetworkHandle,
    0                   // Pinned to Core 0 (Protocol Core)
);
\end{lstlisting}

\section{Bảo vệ vùng nhớ dùng chung bằng dataMutex (main.cpp)}
Trích đoạn \texttt{TaskSensorsCode} thực hiện khóa ghi dữ liệu để cập nhật bản ghi dùng chung một cách phi xung đột:

\begin{lstlisting}[language=C++]
// Read values from all sensors (Sensors.cpp)
SensorRecord record = sensors.readAll();

// Take dataMutex to safely update global sharedRecord
if (xSemaphoreTake(dataMutex, portMAX\_DELAY) == pdTRUE) {
    sharedRecord.temperature = record.temperature;
    sharedRecord.humidity = record.humidity;
    sharedRecord.mq135\_ppm = record.mq135\_ppm;
    sharedRecord.mq135\_v = record.mq135\_v;
    sharedRecord.dust\_density = record.dust\_density;
    sharedRecord.status = record.status;
    sharedRecord.warning = record.warning;
    
    if (record.status == SENSOR\_OK) {
        hasValidReadings = true;
    }
    strncpy(globalStatus, "Measuring...", sizeof(globalStatus));
    xSemaphoreGive(dataMutex); // Release lock immediately after writing
}
\end{lstlisting}

\section{Nhịp xung kích hồng ngoại đọc cảm biến bụi GP2Y (Sensors.cpp)}
Đoạn mã đọc điện áp Analog cảm biến bụi Sharp GP2Y1010AU0F tuân thủ nghiêm ngặt chu kỳ xung $0.32\text{ms}$ của nhà sản xuất:

\begin{lstlisting}[language=C++]
float Sensors::readGP2Y() {
    // Turn on infrared LED (IRED) - active LOW
    digitalWrite(GP2Y\_LED\_PIN, LOW);
    
    // Wait for 280 microseconds according to sensor specs
    delayMicroseconds(280);
    
    // Read ADC value from GP2Y Analog Input pin
    int rawADC = analogRead(GP2Y\_PIN);
    
    // Wait for 40 microseconds before turning off LED
    delayMicroseconds(40);
    digitalWrite(GP2Y\_LED\_PIN, HIGH);
    
    // Convert to actual voltage at 3.3V logic level
    float voltage = (rawADC / 4095.0) * 3.3;
    
    // Calculate dust density using slope and calibration offset
    float dustDensity = (0.17 * voltage - GP2Y\_ZERO\_DUST\_VOLTAGE) * 1000.0;
    
    return (dustDensity < 0.0) ? 0.0 : dustDensity;
}
\end{lstlisting}

\section{Đóng gói JSON và HTTP POST truyền dữ liệu Supabase (Cloud.cpp)}
Đoạn mã serialize dữ liệu và truyền qua HTTPS, xử lý các trường hợp cảm biến lỗi bằng giá trị \texttt{nullptr}:

\begin{lstlisting}[language=C++]
bool Cloud::publishToSupabase(const SensorRecord& data) {
    HTTPClient http;
    if (!http.begin(SUPABASE\_URL)) {
        return false;
    }

    // Set REST API Headers required by Supabase
    http.addHeader("Content-Type", "application/json");
    http.addHeader("apikey", SUPABASE\_KEY);
    String authHeader = "Bearer " + String(SUPABASE\_KEY);
    http.addHeader("Authorization", authHeader.c\_str());

    // Create JSON Document
    StaticJsonDocument<256> doc;
    
    // If sensor value is NaN, map to nullptr to store null in DB
    if (isnan(data.temperature)) doc["temperature"] = nullptr;
    else doc["temperature"] = data.temperature;

    if (isnan(data.humidity)) doc["humidity"] = nullptr;
    else doc["humidity"] = data.humidity;

    if (isnan(data.mq135\_ppm)) doc["mq135\_ppm"] = nullptr;
    else doc["mq135\_ppm"] = data.mq135\_ppm;

    if (isnan(data.dust\_density)) doc["dust\_density"] = nullptr;
    else doc["dust\_density"] = data.dust\_density;

    doc["status"] = (int)data.status;
    doc["warning"] = data.warning;

    String jsonPayload;
    serializeJson(doc, jsonPayload);

    // Send HTTP POST payload
    int httpResponseCode = http.POST(jsonPayload);
    http.end();

    return (httpResponseCode == 200 || httpResponseCode == 201);
}
\end{lstlisting}

\section{Cơ chế ghi cache offline vào bộ nhớ flash (DataManager.cpp)}
Đoạn mã ghi dữ liệu lỗi vào file \texttt{/cache.txt} thông qua hệ thống tệp tin SPIFFS:

\begin{lstlisting}[language=C++]
bool DataManager::cacheRecord(const SensorRecord& record) {
    // Open file in append mode to append to the end
    File file = SPIFFS.open(\_cachePath, FILE\_APPEND);
    if (!file) {
        return false;
    }

    StaticJsonDocument<256> doc;
    doc["temperature"] = isnan(record.temperature) ? nullptr : record.temperature;
    doc["humidity"]    = isnan(record.humidity) ? nullptr : record.humidity;
    doc["mq135\_ppm"]   = isnan(record.mq135\_ppm) ? nullptr : record.mq135\_ppm;
    doc["dust\_density"] = isnan(record.dust\_density) ? nullptr : record.dust\_density;
    doc["status"]      = (int)record.status;

    String jsonString;
    serializeJson(doc, jsonString);
    
    file.println(jsonString); // Write to file and append newline
    file.close();
    return true;
}
\end{lstlisting}
